% Flowchart set: packaging and installation, current state and improved.
% Compile with: xelatex run_installer_flowchart.tex
\documentclass[a4paper,11pt]{article}

\usepackage[margin=18mm]{geometry}
\usepackage{xeCJK}
\usepackage{tikz}
\usetikzlibrary{arrows.meta, positioning}

% CJK fonts -- adjust to whatever is installed on the build machine.
% Common alternatives: "Source Han Sans SC", "WenQuanYi Micro Hei", "SimSun".
\setCJKmainfont{Noto Sans CJK SC}
\setCJKsansfont{Noto Sans CJK SC}
\setCJKmonofont{Noto Sans Mono CJK SC}

\pagestyle{empty}

% Palette: normal steps, steps that touch the system, exits, completion.
\definecolor{stepfill}{HTML}{F1EFE8}
\definecolor{stepline}{HTML}{5F5E5A}
\definecolor{syscallfill}{HTML}{EEEDFE}
\definecolor{syscallline}{HTML}{534AB7}
\definecolor{errfill}{HTML}{FCEBEB}
\definecolor{errline}{HTML}{A32D2D}
\definecolor{donefill}{HTML}{E1F5EE}
\definecolor{doneline}{HTML}{0F6E56}
\definecolor{checkfill}{HTML}{E6F1FB}
\definecolor{checkline}{HTML}{185FA5}

\tikzset{
  % Vertical gap between steps / horizontal gap to the exit branches.
  node distance=9mm and 14mm,
  boxbase/.style={
    rectangle, rounded corners=2pt,
    draw, line width=0.4pt,
    text width=54mm, align=center,
    inner sep=2.6mm, minimum height=15mm
  },
  flowstep/.style={boxbase, fill=stepfill, draw=stepline},
  flowsys/.style={boxbase, fill=syscallfill, draw=syscallline},
  flowdone/.style={boxbase, fill=donefill, draw=doneline},
  flowexit/.style={
    boxbase, fill=errfill, draw=errline,
    text width=42mm, minimum height=11mm
  },
  flow/.style={-{Stealth[length=2.2mm,width=1.8mm]}, line width=0.5pt, draw=stepline},
  exitflow/.style={flow, draw=errline, dashed},
  title/.style={font=\bfseries},
  sub/.style={font=\small},
  % --- compact variants for the 16-step serpentine on page 2 ---
  sboxbase/.style={
    rectangle, rounded corners=2pt,
    draw, line width=0.4pt,
    text width=42mm, align=center,
    inner sep=2mm, minimum height=17mm
  },
  sstep/.style={sboxbase, fill=stepfill, draw=stepline},
  ssys/.style={sboxbase, fill=syscallfill, draw=syscallline},
  scheck/.style={sboxbase, fill=checkfill, draw=checkline},
  sroll/.style={sboxbase, fill=errfill, draw=errline},
  sdone/.style={sboxbase, fill=donefill, draw=doneline},
  soptional/.style={dash pattern=on 1.6pt off 1.2pt},
  rollflow/.style={flow, draw=errline, dashed}
}

% Two-line box content: bold title on top, small subtitle underneath.
\newcommand{\boxtext}[2]{{\bfseries #1}\\[1.4mm]{\small #2}}

% Numbered three-line box content for the serpentine chart.
\newcommand{\sboxtext}[3]{{\bfseries #1.\,#2}\\[1mm]{\footnotesize #3}}

\begin{document}

\begin{center}

{\Large\bfseries 现有打包流程（构建机侧）}

\vspace{5mm}

% Same serpentine layout as page 2: three columns, alternating direction.
\begin{tikzpicture}[node distance=10mm and 10mm]

  % --- row 1: left to right ---
  \node[sstep] (p1) {\sboxtext{1}{准备工作区}{临时目录与产物路径}};
  \node[sstep, right=of p1] (p2) {\sboxtext{2}{校验入参}{检查源码目录存在}};
  \node[sstep, right=of p2] (p3) {\sboxtext{3}{清理源码目录}{执行 make clean}};

  % --- row 2: right to left ---
  \node[sstep, below=of p3] (p4) {\sboxtext{4}{采集版本信息}{四项，含回退值}};
  \node[sstep, left=of p4] (p5) {\sboxtext{5}{收集源文件清单}{递归查找 \texttt{.c} 与 \texttt{.h}}};
  \node[ssys, left=of p5] (p6) {\sboxtext{6}{生成 Makefile}{版本改读固化值}};

  % --- row 3: left to right ---
  \node[ssys, below=of p6] (p7) {\sboxtext{7}{打包源码}{扁平化压缩}};
  \node[ssys, right=of p7] (p8) {\sboxtext{8}{生成安装脚本段}{末尾写分隔标记}};
  \node[ssys, right=of p8] (p9) {\sboxtext{9}{拼接与赋权}{脚本段 + 压缩包}};

  % --- row 4 ---
  \node[sdone, below=of p9] (p10) {\sboxtext{10}{输出产物}{打印 \texttt{.run} 路径}};

  % --- serpentine connections ---
  \draw[flow] (p1) -- (p2);
  \draw[flow] (p2) -- (p3);
  \draw[flow] (p3) -- (p4);
  \draw[flow] (p4) -- (p5);
  \draw[flow] (p5) -- (p6);
  \draw[flow] (p6) -- (p7);
  \draw[flow] (p7) -- (p8);
  \draw[flow] (p8) -- (p9);
  \draw[flow] (p9) -- (p10);

\end{tikzpicture}

\vspace{4mm}

{\small 灰色为准备与采集；紫色为产出中间件；绿色为交付产物。\\
注：步骤 3 在驱动源码目录内执行，会改动被打包的工作区；\\
步骤 5 为递归收集，子目录下的同名文件在步骤 7 扁平化后会互相覆盖。}

\end{center}
\newpage
\begin{center}

{\Large\bfseries 改进后的打包流程（构建机侧）}

\vspace{5mm}

% Same serpentine layout as the installer pages: three columns, alternating
% direction, every node placed relative to the previous one.
\begin{tikzpicture}[node distance=10mm and 10mm]

  % --- row 1: left to right ---
  \node[sstep] (q0) {\sboxtext{0}{入口分派}{参数解析与 \texttt{-{}-help}}};
  \node[sstep, right=of q0] (q1) {\sboxtext{1}{校验入参}{目录、仓库、必需文件}};
  \node[sstep, right=of q1] (q2) {\sboxtext{2}{准备工作区}{退出时自动清理}};

  % --- row 2: right to left ---
  \node[sstep, below=of q2] (q3) {\sboxtext{3}{采集版本信息}{四项 + 打包时间}};
  \node[scheck, left=of q3] (q4) {\sboxtext{4}{工作区状态检查}{未提交改动告警}};
  \node[sstep, left=of q4] (q5) {\sboxtext{5}{收集源文件清单}{白名单、拒绝重名}};

  % --- row 3: left to right ---
  \node[sstep, below=of q5] (q6) {\sboxtext{6}{导出源码副本}{不改动源码目录}};
  \node[ssys, right=of q6] (q7) {\sboxtext{7}{生成 Makefile}{版本改读固化值}};
  \node[scheck, right=of q7] (q8) {\sboxtext{8}{构建配置比对}{与原始 Makefile 核对}};

  % --- row 4: right to left ---
  \node[ssys, below=of q8] (q9) {\sboxtext{9}{打包源码}{生成 \texttt{.tar.gz}}};
  \node[ssys, soptional, left=of q9] (q10) {\sboxtext{10}{预编译模块矩阵}{可选，生成清单}};
  \node[sstep, left=of q10] (q11) {\sboxtext{11}{计算压缩包校验和}{拼接前先算好}};

  % --- row 5: left to right ---
  \node[ssys, below=of q11] (q12) {\sboxtext{12}{注入脚本模板}{一次替换全部占位符}};
  \node[scheck, right=of q12] (q13) {\sboxtext{13}{脚本段自检}{语法与占位符残留}};
  \node[ssys, line width=0.9pt, right=of q13] (q14) {\sboxtext{14}{拼接与赋权}{最后一个写操作}};

  % --- row 6 ---
  \node[sdone, below=of q14] (q15) {\sboxtext{15}{只读自验证与报告}{不再修改产物}};

  % --- serpentine connections ---
  \draw[flow] (q0)  -- (q1);
  \draw[flow] (q1)  -- (q2);
  \draw[flow] (q2)  -- (q3);
  \draw[flow] (q3)  -- (q4);
  \draw[flow] (q4)  -- (q5);
  \draw[flow] (q5)  -- (q6);
  \draw[flow] (q6)  -- (q7);
  \draw[flow] (q7)  -- (q8);
  \draw[flow] (q8)  -- (q9);
  \draw[flow] (q9)  -- (q10);
  \draw[flow] (q10) -- (q11);
  \draw[flow] (q11) -- (q12);
  \draw[flow] (q12) -- (q13);
  \draw[flow] (q13) -- (q14);
  \draw[flow] (q14) -- (q15);

\end{tikzpicture}

\vspace{4mm}

{\small 灰色为准备与采集；紫色为产出文件；蓝色为校验类步骤；\\
虚线边框为可选步骤；绿色为交付产物；粗框的步骤 14 是最后一个写操作。\\
步骤 6 起全部操作在副本上进行，打包对驱动源码目录只读；\\
步骤 14 之后成品不再被修改，步骤 15 只做只读校验。}

\end{center}
\newpage

\begin{center}

{\Large\bfseries 现有安装流程（目标机侧）}

\vspace{6mm}

\begin{tikzpicture}

  % --- main path, top to bottom ---
  \node[flowstep] (env) {\boxtext{环境自检}{工具链、内核头文件}};

  \node[flowstep, below=of env] (work)
    {\boxtext{建立工作目录}{\texttt{/tmp} 下随机命名}};

  \node[flowstep, below=of work] (extract)
    {\boxtext{自解压源码}{跳过脚本段取压缩包}};

  \node[flowstep, below=of extract] (build)
    {\boxtext{编译驱动}{默认 Release 版}};

  \node[flowsys, below=of build] (install)
    {\boxtext{安装模块}{清旧 \texttt{.ko} 后拷入}};

  \node[flowsys, below=of install] (depmod)
    {\boxtext{更新模块索引}{建立硬件到模块的对应}};

  \node[flowdone, below=of depmod] (finish)
    {\boxtext{清理并退出}{删除工作目录}};

  % --- exit branches, placed relative to the step they leave from ---
  \node[flowexit, right=of env] (exitdep) {缺依赖则退出};

  \node[flowexit, right=of build] (exitbuild) {编译失败退出\\[0.8mm]{\small 源码残留}};

  % --- connections ---
  \draw[flow] (env)     -- (work);
  \draw[flow] (work)    -- (extract);
  \draw[flow] (extract) -- (build);
  \draw[flow] (build)   -- (install);
  \draw[flow] (install) -- (depmod);
  \draw[flow] (depmod)  -- (finish);

  \draw[exitflow] (env)   -- (exitdep);
  \draw[exitflow] (build) -- (exitbuild);

\end{tikzpicture}

\vspace{4mm}

{\small 紫色为改动系统状态的步骤；虚线为异常退出路径。}

\end{center}

\newpage
\begin{center}

{\Large\bfseries 改进后的安装流程（目标机侧）}

\vspace{5mm}

% Serpentine (boustrophedon) layout: three columns, alternating direction,
% so all 16 steps fit on one A4 page. Every node is placed relative to the
% previous one -- no absolute coordinates.
\begin{tikzpicture}[node distance=10mm and 10mm]

  % --- row 1: left to right ---
  \node[sstep] (s0) {\sboxtext{0}{入口分派}{参数与 \texttt{-{}-version}}};
  \node[sstep, right=of s0] (s1) {\sboxtext{1}{权限自检}{统一检查 root}};
  \node[sstep, right=of s1] (s2) {\sboxtext{2}{环境自检}{工具链、内核树、余量}};

  % --- row 2: right to left ---
  \node[sstep, below=of s2] (s3) {\sboxtext{3}{安装包自校验}{校验源码包完整性}};
  \node[sstep, left=of s3] (s4) {\sboxtext{4}{建立工作目录}{内存盘优先、限额}};
  \node[sstep, left=of s4] (s5) {\sboxtext{5}{解压源码}{展开到工作目录}};

  % --- row 3: left to right ---
  \node[sstep, soptional, below=of s5] (s6) {\sboxtext{6}{预编译匹配}{命中则跳过编译}};
  \node[sstep, right=of s6] (s7) {\sboxtext{7}{编译驱动}{默认 Release 版}};
  \node[ssys, right=of s7] (s8) {\sboxtext{8}{备份现有模块}{供回滚使用}};

  % --- row 4: right to left ---
  \node[ssys, below=of s8] (s9) {\sboxtext{9}{安装模块}{清旧后拷入}};
  \node[ssys, left=of s9] (s10) {\sboxtext{10}{更新模块索引}{建立硬件对应}};
  \node[scheck, left=of s10] (s11) {\sboxtext{11}{安装后验证}{试加载并确认绑定}};

  % --- row 5: left to right ---
  \node[ssys, soptional, below=of s11] (s12) {\sboxtext{12}{配置节点权限}{可选}};
  \node[sroll, right=of s12] (s13) {\sboxtext{13}{失败回滚}{恢复备份模块}};
  \node[sdone, right=of s13] (s14) {\sboxtext{14}{统一清理}{覆盖所有退出路径}};

  % --- row 6 ---
  \node[sdone, below=of s14] (s15) {\sboxtext{15}{结果汇报}{日志落至 \texttt{/tmp}}};

  % --- serpentine connections ---
  \draw[flow] (s0)  -- (s1);
  \draw[flow] (s1)  -- (s2);
  \draw[flow] (s2)  -- (s3);
  \draw[flow] (s3)  -- (s4);
  \draw[flow] (s4)  -- (s5);
  \draw[flow] (s5)  -- (s6);
  \draw[flow] (s6)  -- (s7);
  \draw[flow] (s7)  -- (s8);
  \draw[flow] (s8)  -- (s9);
  \draw[flow] (s9)  -- (s10);
  \draw[flow] (s10) -- (s11);
  \draw[flow] (s11) -- (s12);
  % Normal path skips the rollback node: detour below row 5 so the line
  % never crosses node 13, and enter node 14 left of its exit to node 15.
  \draw[flow] (s12.south) -- ++(0,-7mm) -| ([xshift=-13mm]s14.south);
  \draw[flow] ([xshift=13mm]s14.south) -- ([xshift=13mm]s15.north);

  % Rollback is not part of the normal order: any failing step enters it,
  % and it merges back into the common cleanup.
  \draw[rollflow] (s12) -- (s13);
  \draw[rollflow] (s13) -- (s14);

\end{tikzpicture}

\vspace{4mm}

{\small 灰色为准备与构建；紫色为改动系统状态；蓝色为安装后验证；\\
虚线边框为可选步骤；红色虚线为异常回滚路径，任一步失败后进入 13，再汇入 14。}

\end{center}

\end{document}
